\section{Exact certification and adversarial controls}

The executable package uses rational arithmetic throughout.  The producer checks every
$Q=8,\ldots,512$ and every $L=1,\ldots,4$.  Across these 505 scales, the first
three dilations have no primitive collision.  The fourth dilation has 182
collision-bearing scales and 235 resonance pairs; the first occurs at $Q=25$.
These counts certify the finite census only.  They are not a density theorem for the
linear prime relation $7p+3r=16Q$.

Thirty additional records evaluate aligned, affine, and balanced sign profiles at ten
collision-bearing scales through $Q=1000$.  Each record rebuilds the literal rows,
their supports, all four energies, and the signed correction independently of any
floating-point tolerance.  The independent checker does not import the producer and
runs identically under normal and optimized Python modes.

The adversarial test targets the most likely source error.  If primitivity is removed at
$L=3$, $Q=8$, then the primes $11$ and $13$ appear to collide through multipliers
$4$ and $-4$ at two coordinates.  Both multipliers have nontrivial gcd with $h=96$,
so the literal primitive row contains neither point.  Restoring $(m,h)=1$ removes the
entire false transition.  The same adversary then checks the valid $Q=25$ residues and
the three energy signs from Table~\ref{tab:profiles}.

The certificate and checker are trust boundaries, not theorem evidence.  The theorem is
the integer classification in Section~3; the computation tests implementation,
boundary, schema, and exact-energy consistency.
